Trước hết, tôi xin bày tỏ lòng biết ơn sâu sắc tới TS. Ma Thị Châu vì sự hỗ trợ và chỉ dẫn tận tâm trong suốt quá trình thực hiện khóa luận này. Những định hướng quý báu, chuyên môn sâu rộng, phản hồi mang tính xây dựng và sự tận tụy của các thầy cô đã đóng vai trò then chốt đối với thành công của tôi. Tôi trân trọng cơ hội được học hỏi và làm việc dưới sự hướng dẫn của cô.

Tôi cũng xin gửi lời cảm ơn đến các cộng sự nghiên cứu của tôi gồm Nguyễn Ngọc Hưng và Nguyễn Minh Kiên, cùng người bạn thân thiết Phạm Thu Trang. Sự hỗ trợ và động viên của họ là vô cùng quý giá trong suốt bốn năm học đại học của tôi.

Những kết quả trong luận văn này đóng góp vào đề tài QG.25.03 “Nghiên cứu tái tạo chuyển động diễn xuất khuôn mặt nhân vật trên mô hình 3D, phục vụ công tác giảng dạy sân khấu truyền thống”.

Ngoài ra, tôi xin gửi lời cảm ơn đến sự hợp tác nhiệt tình của PGS.TS Đinh Quang Trung, Viện trưởng Viện Sân khấu - Điện ảnh, trường Đại học Sân khấu - Điện ảnh Hà Nội, cùng cô Hoàng Thị Trần Chuyến và các bạn học sinh của cô tại Trường THPT Việt Nam - Ba Lan. Chuyên môn và sự tận tụy của họ đã đóng vai trò quan trọng trong việc hỗ trợ quá trình thử nghiệm thực tế và hoàn thiện liên tục các kết quả nghiên cứu của tôi, đảm bảo tính ứng dụng cao cho công trình này.
